\documentclass[10pt]{article}
\usepackage[margin=20mm]{geometry}
\usepackage{amsmath,amssymb,mathtools,fontspec,unicode-math,xurl,hyperref}
\IfFontExistsTF{DejaVu Serif}{\setmainfont{DejaVu Serif}}{\setmainfont{Latin Modern Roman}}
\setmathfont{latinmodern-math.otf}
\setlength{\parindent}{0pt}\setlength{\parskip}{6pt}\sloppy
\title{Scientific initial review of
Abelian differentials and their periods:
The bi-algebraic point of view}
\author{AI-assisted technical review}
\date{26 October 2026}
\begin{document}
\maketitle
Abstract: The assigned arXiv v1 is reviewed through its arithmetic characterization, distribution argument and conditional linearity proofs. Six local correction groups concern quotient conventions, the minimal abelian factor, the square-tiled terminology, a projection asserted to be an isomorphism, and two stages of the linearity argument. Explicit examples and bounded repairs are supplied. In particular, a weaker dimension inequality preserves the distribution contradiction without the asserted isomorphism. Later published nonlinear examples contradict Conjecture 2.7, rather than the conditional linearity theorems. The assessment is Minor revision for this exact author version. External foundational results are not all independently reproved, and scientific adjudication and human review remain pending.\par
\section*{Identity, version and scope}
The authors are Bruno Klingler and Leonardo A. Lerer. The journal association is DOI https://doi.org/10.4171/JEMS/1581. The actual reviewed input is the 28-page arXiv:2202.06031v1, uploaded on 12 February 2022; its PDF and printed page numbers agree. EMS records submission on 22 December 2022, acceptance on 26 February 2024 and online-first publication on 3 January 2025. The publisher PDF is subscription-restricted and was not the reviewed input. The separately supplied author-hosted Ax-AbelianDifferentialsv8.pdf has different bytes and no explicit version declaration. Its file creation metadata is not an acceptance date. Neither its equivalence with v1 nor equivalence of either author file with the Version of Record has been established.\par
The arithmetic condition uses \(\overline{\mathbb Q}\), not \(\mathbb Q\): both the algebraic differential up to scaling and its projective relative-period vector must be algebraic. The distinction is essential. For genus at least two and \(n\) zeros, the relative-cohomology rank is \(2g+n-1\), whereas the projectivized stratum has dimension \(2g+n-2\). The level structure and the finite cover labeling zeros must be respected when using local systems.\par
Theorem 2.2 involves the smallest abelian quotient \(A_{[\omega]}\) carrying the differential, its isotypic CM property, a central-CM-field eigenline, and torsion differences between images of zeros. Theorem 2.6 is a nondensity statement for invariant linear loci with the specified rank and field degree; it does not assert nondensity in every stratum. Theorem 2.8 concerns curves satisfying \((\star)\), while Theorem 2.9 concerns subvarieties in isoperiodic leaves. The constant-part case in which the projective differential varies is expressly left outside those two arguments. Conjecture 2.7 is therefore a separate assertion.\par
\section*{Established local problems}
Problem 1 (Minor: the quotient and projective-dimension conventions need correction).  Location. Section 1.2.2, p. 3; Section 3.5(i), p. 9; compare Section 6.1, p. 20.\par
The foliation on \(S_\alpha\) is described on p. 3 as having codimension \(2g\). In period coordinates it is given by \[\mathbb P V_\alpha\dashrightarrow\mathbb P V_\alpha^p,
 \qquad \dim\mathbb P V_\alpha^p=2g-1.\] The map is a submersion where the pure component is nonzero. Consequently the codimension is \(2g-1\) on \(S_\alpha\), and \(2g\) on \(H_\alpha\). For example, the minimal genus-two stratum \(S_2\) has dimension three and zero-dimensional isoperiodic leaves; codimension four is impossible. Likewise, a projective chart on an \(m\)-dimensional manifold requires \(\operatorname{rank}V_\mathbb Z=m+1\), not \(m\) as printed on p. 9.\par
There is also an omitted quotient in the Hurwitz-space discussion. With the stated meaning of covers to a fixed elliptic target, fix \(f:C\to E\) with \(g(C)\ge2\). The family \(t_a\circ f\), \(a\in E\), consists generically of distinct covers, since \(\operatorname{Aut}(C)\) is finite. Yet \[(t_a\circ f)^*\omega_E=f^*\omega_E\] for every \(a\). Thus the displayed map to projectivized differentials has a positive-dimensional fiber and cannot be an immersion.\par
Impact and repair boundary. Distinguish the two foliations and correct the projective rank. Define the Hurwitz space modulo target translations, or explicitly work with its image and an appropriate normalization; the unquotiented immersion assertion must be removed. These corrections preserve the intended isoperiodic image. The genus-one case should also be stated separately: the formulas \(4g-4\) and \(2g+n-2\) on p. 7, with positive zero orders and hence \(n=0\) at \(g=1\), do not describe the one-dimensional modular curve \(S_0\) used later. The elliptic case already has its own argument on pp. 12 and 14.\par
Problem 2 (Minor: the smallest factor is replaced by an entire isotypic block).  Location. Definition 2.1, p. 4; the identification \(A_{[\omega]}:=A_\nu^{k_\nu}\) in the proof of Theorem 2.2, p. 13.\par
Consider the smooth genus-two curve and elliptic curve \[C:\ y^2=x^6-1,\qquad E:\ Y^2=X^3-1.\] There are two maps, extending over the smooth projective models, \[f(x,y)=(x^2,y),\qquad
 h(x,y)=(x^{-2},iy/x^3).\] Their differential pullbacks are \[f^*(dX/Y)=2x\,dx/y,\qquad h^*(dX/Y)=2i\,dx/y.
 \tag{2.1}\] They are independent, so the induced map \(\operatorname{Jac}(C)\to E^2\) is an isogeny. For \(\omega=2x\,dx/y\), however, the minimal factor is \(A_{[\omega]}\sim E\), of dimension one. The full nonzero isotypic block of \(\operatorname{Jac}(C)\) is \(E^2\), of dimension two. The printed identification is therefore false.\par
Impact and repair boundary. Decompose the minimal quotient \(A_{[\omega]}\) itself, or retain the distinction between that quotient and the larger ambient isotypic block. Minimality, together with the period argument, then gives the required isotypy. The degree must remain the dimension of the minimal quotient. This example satisfies the intended CM characterization; it contradicts the intermediate identification, not Theorem 2.2.\par
Problem 3 (Minor: general CM elliptic covers are not all square-tiled in the stated sense).  Location. Remark 5.1(1) and the final paragraph of Section 5.1, pp. 15–16.\par
Use \(C,E,f,\omega\) from Problem 2. The degree-two map \(f\) ramifies only at the two points with \(x=0\); their images \((0,\pm i)\) on \(E\) are nonzero three-torsion points. Indeed the tangent lines \(Y=\pm i\) meet \(E\) to order three there. Hence the construction of Section 5.1 gives an arithmetic point of degree one.\par
It cannot, even after complex rescaling, have all its periods in \(\mathbb Q(i)\), as required by Remark 5.1(1). The rational span of its absolute periods is \(\Omega K\), where \(K=\mathbb Q(\sqrt{-3})\): the image of \(f_*H_1(C,\mathbb Z)\) has finite index in \(H_1(E,\mathbb Z)\). If a rescaling put this two-dimensional rational space inside \(\mathbb Q(i)\), it would equal \(\mathbb Q(i)\). The ratio of two independent periods would then force \[\mathbb Q(\sqrt{-3})=\mathbb Q(i),\] which is impossible.\par
Impact and repair boundary. Composing with \([N]:E\to E\) does collapse the finite torsion branch locus to the origin, and multiplies the differential by \(N\). It does not change the CM field or turn \(E\) into a square torus. Replace the last claim by the correct statement about covers of the same CM elliptic curve branched over one point; restrict to the \(\mathbb Q(i)\) case when using the stated square-tiled terminology. The density proof using \(\mathbb Q(i)\) remains valid, as does the degree-one arithmetic construction.\par
Problem 4 (Minor: Lemma 5.10 asserts an unnecessary isomorphism).  Location. Lemma 5.10, equations (5.9)–(5.12), p. 19.\par
The projection to the distinguished abelian factor need not recover the covering curve. A concrete test is the genus-two, three-square origami in \(H(2)\). Connected degree-three covers of a once-punctured torus with a three-cycle commutator are described by pairs \((h,v)\in S_3^2\), up to simultaneous conjugacy. There are exactly three classes, with cycle-order types \[A=(2,2),\qquad B=(2,3),\qquad C=(3,2).\] An independent finite enumeration found 18 pairs and these three conjugacy classes. The Nielsen operations \[S(h,v)=(v^{-1},h),\qquad T(h,v)=(h,vh)\] act as \(S=(BC)\) and \(T=(AB)\). They therefore belong to one Teichmüller curve and give a generically three-to-one map to the elliptic modular curve. Riemann–Hurwitz gives genus two and one zero of order two. Here \(k=d=1\), \(u=0\), \(n=1\), so the target in (5.8) is precisely the elliptic factor, without relative marked-point data.\par
This multiplicity is not automatically removed by the permitted level convention. For level five, \(T^5\) and \(ST^5S^{-1}\) lie in \(\Gamma(5)\) and still act transitively on the three classes. Adding level data to the genus-two curve refines a covering; it cannot identify these distinct generic differentials over the same elliptic data. The assertion of isomorphism is thus too strong. More generally, the claim that the whole mixed extension is recovered by the displayed projection needs an additional reconstruction argument; forgetting a complementary abelian factor does not supply one.\par
Impact and repair boundary. The distribution argument can avoid this assertion entirely. Let \[D=d\left(\frac{k(k+1)}2+uk\right).\] Assuming the cited special-locus and monodromy computation (5.11), the automatic inequality \(\dim\overline{Z'}=\dim Z'\le\dim Z=u+2k-1\) suffices. For \(d\ge2\), \[D-(u+2k-1)\ge
 \begin{cases}
 (d-1)(u+1)>0,&k=1,\\
 k^2-k+1+(2k-1)u>0,&k\ge2.
 \end{cases}
 \tag{2.2}\] This is the required contradiction. Delete Lemma 5.10 or replace it by the weaker statement actually needed, and replace (5.12) by the inequality. The example does not refute Theorem 2.6, and no independent reproof of every cited monodromy theorem is claimed.\par
Problem 5 (Minor: the Ax–Schanuel iteration restricts the wrong inverse image).  Location. Proof of Proposition 6.3, p. 21, after the first application of Theorem 6.4.\par
From \(\sigma(Z)\subset a+B\), the proof infers that \(W=\operatorname{pr}^{-1}(V)\) is contained in a translate of \(L=\operatorname{Lie}B\). Only the lifted image of \(Z\) is so contained. The full inverse image \(W\) also contains all directions in \(\ker\operatorname{pr}\). Already for a point \(Z\), with \(r\ge2\), \(n=2\) and \(B=0\), this inverse image has dimension \(r-1\), whereas \(L=0\). The asserted containment does not follow even in this degenerate case allowed by the proposition.\par
Impact and repair boundary. Translate the lifted image into \(L\), and at the next stage use \[W_L=(\operatorname{pr}|_L)^{-1}(V)\subset L.\] If \(V\subsetneq\operatorname{pr}(L)\), then \(\dim W_L\le\dim L-1\). With \(N'=\dim L\) and \(s=\dim\overline{\sigma(Z)}\), the same computation gives \[\operatorname{codim}(\overline{\sigma(Z)}\times W_L)\ge N'-s+1,\quad
 \operatorname{codim}\Delta_B=N',\quad
 \operatorname{codim}U\le2N'-s.\] Ax–Schanuel again lowers the abelian dimension. Iteration ends with \(V\) equal to an affine image of a linear space. Bi-algebraicity supplies the required dimension equality between the original local image and its algebraic model. This fixes the written iteration without asserting a counterexample to Proposition 6.3 or Theorem 2.9.\par
Problem 6 (Minor: omitted hypotheses in the final linearity steps).  Location. Theorem 2.8, Step 2, p. 24; Lemma 6.7(2), p. 24; Proposition 6.9, pp. 25–26.\par
First, the existence of a dense orbit in \(\mathcal Y\) does not by itself imply that its linear span \(W\) is an irreducible representation. For instance, under \(\mathrm{SL}_2\) on \(\mathbb C^2\oplus\mathbb C\), the affine plane \(\mathbb C^2\times\{1\}\) has a dense orbit and spans the reducible representation \(\mathbb C^2\oplus\mathbb C\).\par
The actual argument has sufficient additional structure. Let \(K=\ker(W\to\mathcal Y^p)\subset W_0V_\alpha\). The connected semisimple group \(H'\) acts trivially on \(K\), since zero-label monodromy is finite. Complete reducibility gives \(W\simeq\mathcal Y^p\oplus K\). The orbit of a point \((v,k)\) with \(v\ne0\) is \((\mathcal Y^p\setminus\{0\})\times\{k\}\); its two-dimensional closure is all of the irreducible \(\mathcal Y\). Since \(\mathcal Y\) is a cone, \(k=0\). Since it spans \(W\), \(K=0\), proving the intended conclusion. These missing steps, rather than the false general inference, should be written.\par
Second, Lemma 6.7(2) uses \(\dim\widetilde{\mathrm{SO}}_{2n+1}\) although its parabolic belongs to \(\widetilde{\mathrm{SO}}_{2n}\). At \(n=3\), the printed difference is \(21-12=9\), whereas the claimed binomial value is three. The correct difference is \(15-12=3\). Use the even orthogonal group and consistently count projective highest-weight orbits in the preceding exclusions.\par
Third, the appeal to quasi-finiteness in Proposition 6.9 is made for the open set \(U\) of nonzero forms. Its period map forgetting scale has \(\mathbb C^*\)-fibers. Work on \(U/\mathbb C^*\), or use the direct genus-two monodromy argument: the cover labeling the two zeros, modulo scaling, is \(C\) minus its Weierstrass points. Moving one zero along a nonzero homology cycle \(c\) moves the other along its hyperelliptic conjugate, and changes a connecting relative path by \[\iota_*c-c=-2c\ne0.\] For generic \(\omega\), this yields the infinite period coset used on p. 26 and contradicts generic finiteness of a hypothetical two-dimensional algebraic graph.\par
Impact and repair boundary. Each correction supplies a local step using the existing hypotheses. The examples refute the printed inference or formula, while the repaired arguments preserve the stated conditional linearity and genus-two conclusions.\par
\section*{Positive checks, later counterevidence and limits}
\subsection*{Arithmetic and distribution mechanisms}
The CM eigenline criterion is compatible with the period decomposition: in the simple CM case, distinct embeddings have independent period lines, so an algebraic differential with a single projective algebraic period line must select one embedding. In an isotypic power the center, rather than a chosen individual copy, is the relevant action. The minimal quotient must be fixed first, as in Problem 2.\par
For the relative-period argument in Lemma 4.3, a relation between a relative period and a nonzero absolute period defines the algebraic hyperplane in \(\operatorname{Lie}(A\times A)\). The stated analytic subgroup theorem produces the requisite algebraic subgroup. In the simple case its nonzero first projection is an isogeny, yielding an endomorphism relation and then membership in the CM period \(\overline{\mathbb Q}\)-structure; the quoted torsion theorem finishes the implication. The full general-isotypic reduction and the imported transcendence theorems were not independently reproved. No false torsion conclusion was established. Conversely, torsion zero differences multiply relative paths into absolute cycles, so the CM eigenline condition supplies algebraic ratios of all relative periods.\par
The \(\mathbb Q(i)\)-coordinate density argument on pp. 14–15 has the correct local structure. Clearing denominators gives integration into a square torus; zeros map to the origin and the branched cover descends over \(\overline{\mathbb Q}\). Projectivization removes the arbitrary complex normalization of the elliptic differential. Problem 3 changes the claimed exhaustive description, not this dense subset. For a finite cover \(f\), the rational cohomology pullback is injective and the minimal Hodge structure carrying \(f^*\omega\) is isomorphic to the original one; this supports the degree preservation in Lemma 5.2.\par
The example \(y^2=x^7-1\), \(\omega=x\,dx/y\), on p. 17 has divisor \((0,i)+(0,-i)+2\infty\). The function \(y-i\) has divisor \(7(0,i)-7\infty\), giving the stated torsion test. The CM simplicity and the non-Veech conclusion retain the cited external dependencies. For the exact-degree step in Theorem 2.6, the totally real field \(K_0\) acts on the minimal quotient carrying its eigendifferential. If that quotient has dimension \(a\), each real embedding of \(K_0\) gives a real summand of \(H_1\) preserved by the complex structure, so \(2a/d\) is even and \(d\mid a\). Together with \(a\le kd\), this gives \(a=d\) when \(k=1\), and \(a=kd\) under the higher-rank degree hypothesis. Arithmetic images and André–Oort then feed the special-locus dimension formula; the strict inequality in (2.2) verifies the final numerical obstruction for every \(d\ge2\), \(k\ge1\), \(u\ge0\).\par
\subsection*{Conditional linearity}
The pure-period model of a curve satisfying \((\star)\) is one-dimensional and has nontrivial semisimple monodromy action. Highest-weight orbit dimensions exclude all the stated weight-one representation types except the standard \(\mathrm{SL}_2\) case. For type \(A_n\), the projective orbit has dimension \(r(n+1-r)\); the spinor varieties have dimensions \(n(n+1)/2\) and \(n(n-1)/2\). For type \(C_n\) the projective orbit has dimension \(2n-1\); for the standard even-orthogonal highest-weight orbit it is \(2n-2\). The low-rank identifications must be treated before these counts. Multiple independent \(\mathrm{SL}_2\) factors or a nonzero trivial component would give an orbit too large for a curve. Problem 6 completes the passage from pure to relative periods.\par
For an isoperiodic algebraic locus the fixed-part theorem provides the constant abelian factor. The corrected Ax–Schanuel iteration makes its relative model affine, and the scaling cone over an affine fiber gives the projective linearity of Theorem 2.9. In genus two, the fixed factor is either elliptic, forcing projectively constant pure periods, or the entire abelian surface. Torelli then reduces the latter case to a fixed curve, where the relative-period monodromy in Problem 6 excludes bi-algebraicity.\par
The quartic \(Y^4=XZ(X-Z)(X-\lambda Z)\), \(\lambda\ne0,1\), in Section 6.4 is smooth of genus three. The vertical pencil has exactly the four stated total tangencies, each of intersection multiplicity four. The two-dimensional constant holomorphic part confines its projective periods to a line. This is consistent with the example and does not settle the unrestricted conjecture.\par
\subsection*{Later evidence against the unrestricted conjecture}
Deroin–Matheus [DM], Theorem 1 and pp. 1693–1696, give \[C_u:\ y^6=x(x-1)(x+1)(x-u),\qquad
 \omega_u=x^2dx/y^5,\qquad u\notin\{0,\pm1\}.\] This is a nonlinear bi-algebraic curve in the minimal genus-seven stratum. Finite monodromy supplies algebraicity. Their displayed Gauss–Manin calculation gives three independent vectors \(p_2,\nabla p_2,\nabla(6(1-u^2)\nabla p_2)\). Recomputing its determinant gives \[-\frac{35}{36(1-u^2)},\qquad
 \left.-\frac{35}{36(1-u^2)}\right|_{u=2}=\frac{35}{108}\ne0.\] The principal checked this last value with exact rational arithmetic. These examples belong to the varying-differential fixed-part regime. They contradict Conjecture 2.7, not the \((\star)\) or isoperiodic theorems. A current account should cite this development. It postdates the reviewed February 2022 version; it is not evidence that the authors claimed a proved unrestricted theorem, nor evidence about uninspected changes in the 2025 publisher version.\par
\subsection*{Coverage, attribution and assessment}
The principal read all 1,737 lines of the supplied primary text extraction and inspected all 28 primary PDF pages. The 511-line extraction and all nine pages, including the cover, of the later Deroin–Matheus publisher PDF were also read and inspected. The principal evaluated and adopted the findings, wrote this report, ran the actual compiler, and inspected every final report page. The auxiliaries are not independent acceptance reviewers or independently authenticated Pro backends.\par
The exact finite \(S_3\) enumeration and rational determinant check were independently written review calculations. No author code or author TeX was executed. Same-version TeX passages were read only for source confirmation. The author-hosted alternative PDF was retained as a distinct input, without whole-version equivalence. A bounded publisher, arXiv, author-page and correction/supplement search identified the later nonlinear paper but no separate correction or executable supplement for this exact v1. Lack of code is not a mathematical defect.\par
The analytic subgroup, CM transcendence, André–Oort, full monodromy, Satake and fixed-part inputs remain external dependencies; reading the manuscript is not an independent reconstruction of all those results. Several auxiliary claims require the explicit corrections above, but the supplied repairs preserve the main routes within that dependency boundary. The assessment is Minor revision for this precise input, with six local problem groups and separately documented later counterevidence to a conjecture. This is an initial review with `scientific\_decision=pending` and `human\_reviews=[]`.\par
9 B. Klingler and L. A. Lerer, Abelian differentials and their periods: the bi-algebraic point of view, arXiv:2202.06031v1, 12 February 2022. https://arxiv.org/abs/2202.06031v1. EMS Press, article record, DOI 10.4171/JEMS/1581. https://ems.press/journals/jems/articles/14298461. B. Deroin and C. Matheus, Non-linear bi-algebraic curves and surfaces in moduli spaces of Abelian differentials, C. R. Math. 361 (2023), 1691–1698. DOI https://doi.org/10.5802/crmath.529. Publisher PDF with repository cover: https://www.numdam.org/item/10.5802/crmath.529.pdf.\par
\end{document}
